\section{Preliminary mixing}\label{pre:mixing}\label{pre:section}

We first prove a lattice estimate whose length scale is less precise than the
natural continuum scale, but whose uniformity under weakening bonds will be
needed later. Spins take values in the unit sphere \(S^3\subset\R^4\), identified
with the unit quaternions. The a priori measure is normalized Haar measure.
For a finite graph in the square lattice the Gibbs density is proportional to
\[
 \exp\left(\sum_{e=xy}\beta_e q_x\cdot q_y\right).
\]
Each unoriented edge occurs once. A free subgraph means that bonds outside the
graph are absent; it imposes no pinned spin values.

For a local Lipschitz observable \(F\), supported on a finite set \(A\), put
\[
 L_F=\norm{F}_\infty+\max_{x\in A}\operatorname{Lip}_x(F),
\]
where the individual Lipschitz constant uses round distance on \(S^3\) with
all other spins fixed. Distances between supports use the maximum norm, or
its periodic version on a torus.

\begin{theorem}[Preliminary mixing]\label{thm:preliminary}
There are fixed positive constants \(C,c,p\) such that, on setting
\begin{equation}
 \Xi(\beta)=\exp\left[\pi\beta+
 C\sqrt{(1+\beta)\log(2+\beta)}\right],
 \label{eq:pre:Xi}
\end{equation}
the following hold for every \(\beta\geq0\).

For every unpinned rectangular torus with each side at least \(C\Xi(\beta)\),
and every choice \(0\leq\beta_e\leq\beta\), local observables \(F,G\) supported
on \(A,A'\) at distance \(d\) satisfy
\begin{equation}
 |\Cov(F,G)|\leq C(1+\beta+|A|+|A'|+d)^{20}
 L_FL_G e^{-cd/\Xi(\beta)}.
 \label{eq:pre:theorem-mixing}
\end{equation}
The same bound holds on every unpinned free subgraph, with ambient lattice
distance. In the homogeneous model the local limit of periodic square
measures is unique.

Let \(Z_\beta(n,m)\) be the partition function of the homogeneous
\(n\)-by-\(m\) torus. For every even \(n\geq C\Xi(\beta)\),
\begin{equation}
 0\leq\Delta_\beta(n):=4\log Z_\beta(n,n)-\log Z_\beta(2n,2n)
 \leq C(1+\beta)^pn^p e^{-cn/\Xi(\beta)}.
 \label{eq:pre:theorem-delta}
\end{equation}
\end{theorem}

The proof first obtains a small stiffness on a suitable pinned square. This
part allows arbitrary signed couplings and fixed right twists. We then use
coarse orientation variables to prove decay for charged observables. Two
conditional decompositions into ferromagnetic circle and Ising systems give
mixing for arbitrary local observables. Only this last part requires
nonnegative untwisted couplings.

Constants denoted by \(C\) may increase from line to line. All are independent
of lattice size and of the particular couplings and pins. Dependence on a
fixed moment order or a fixed auxiliary exponent is indicated when needed.
For continuum test fields on the unit square, \(\|\cdot\|_2\) denotes the
Lebesgue \(L^2\) norm. For random variables, \(\|\cdot\|_p\) denotes the
probabilistic \(L^p\) norm.

\subsection{Local class and Ward identities}\label{pre:local}

The following calculation is a local Ward identity, in the tradition of
\cite{Mermin1967,DriesslerLandauPerez,AizenmanSimonWard}. We give the full
quaternion formula because the noncommuting current term and its factor
two determine the quantitative stiffness gain.

Work on \(Q(S)=\{0,\ldots,S\}^2\) with all boundary spins fixed. For each positively oriented edge \(e=x\to y\), set

\[
 t_e+s_e=q_yU_eq_x^{-1},\qquad t_e\in\mathbb R,\quad s_e\in\operatorname{Im}\mathbb H\simeq\mathbb R^3,
\]

where \(U_e\) is any fixed unit quaternion. Take signed coefficients \(|b_e|\leq b\leq B\), with \(w_e=b_e\) except that edges tangent to the boundary carry \(w_e=b_e/2\). The local density is \(Z^{-1}\exp(\sum_e w_et_e)\). The half allocation makes adjacent cell actions sum exactly to the action on the union.

For macroscopic midpoint coordinates \(z_e\), define

\[
 X_Q(F)=S^{-1}\sum_e w_es_e\cdot F_{\mu(e)}(z_e),
\]

\begin{equation}
 \mathcal H_Q(F,G)=S^{-2}\mathbb E\sum_e w_et_e\,\overline{F_e}\cdot G_e
 -\operatorname{Cov}(X_Q(F),X_Q(G)).
\label{eq:pre:1-1}
\end{equation}

The covariance is sesquilinear, with conjugation on the first entry. Write \(H_Q\) for the restriction to constant \(3\times2\) matrices. In each spatial direction the sum of the absolute edge coefficients is at most \(bS^2\): there are \(S(S-1)\) interior edges and two tangent boundary rows with half weight. Thus

\begin{equation}
H_Q\leq bI.
\label{eq:pre:1-2}
\end{equation}

This is valid with signed coefficients and twists. No positive lower bound on \(H_Q\) is claimed or used.

For an imaginary unit \(a\), rotate \(q_x\mapsto e^{af_x}q_x\). A bond quaternion changes to \(e^{af_y}(t_e+s_e)e^{-af_x}\). Its \((1,a)\)-plane rotates through \(f_y-f_x\); its imaginary plane perpendicular to \(a\) rotates through \(f_y+f_x\). If \(\mathcal A=-\sum w_et_e\), then

\begin{equation}
 D_f\mathcal A=\sum_e w_e(s_e\cdot a)(f_y-f_x),\qquad
 D_f^2\mathcal A=\sum_e w_et_e(f_y-f_x)^2.
\label{eq:pre:1-3}
\end{equation}

For a general infinitesimal imaginary profile \(\eta_x\), the imaginary bond derivative is exactly

\begin{equation}
D_\eta s_e=t_e(\eta_y-\eta_x)+(\eta_y+\eta_x)\times s_e.
\label{eq:pre:1-4}
\end{equation}

Consequently the compact-profile Ward identity has a contact term and a bracket term. If \(\eta\) vanishes on the fixed boundary, Haar integration by parts gives

\begin{equation}
\mathbb E[\overline{X_Q(d_S\eta)}X_Q(F)]
=S^{-2}\mathbb E\sum_e w_et_e\overline{(d_S\eta)_e}\cdot F_e
+\mathbb E X_Q\!\left(F_e\times(\overline{\eta_y}+\overline{\eta_x})\right),
\label{eq:pre:1-5}
\end{equation}

where the cross product is extended complex bilinearly. This follows first
for real profiles and then by complex linear extension. Here
\((d_S\eta)_\mu(z)=S[\eta(z+e_\mu/(2S))-\eta(z-e_\mu/(2S))]\).
This fixes the normalization and the factor two in the eventual nonabelian derivative. For the same-axis score there is no bracket term; \eqref{eq:pre:1-3} follows directly.


\subsection{Current moments near a pinned boundary}\label{pre:moments}

Assume throughout the local induction

\begin{equation}
\log(2+S)\leq16B.
\label{eq:pre:2-1}
\end{equation}

\begin{lemma}[Uniform current moments]\label{pre:current-moments}
For every fixed \(1\leq p<\infty\),

\begin{equation}
\|X_Q(F)\|_p\leq C_p(1+b)\log(2+S)\|F\|_{C^1}.
\label{eq:pre:2-2}
\end{equation}

Here and subsequently the derivative norm is in the indicated rescaled coordinates. The constants are independent of pins, signs, twists, and side length.
\end{lemma}

\begin{proof}

\subsubsection{Compact scores}\label{pre:part-2-1}

For a real one-axis profile \(f\) zero at the pinned boundary, the second variation along the full rotation path is bounded in absolute value by
\(b\sum_e(f_y-f_x)^2\). Taylor's formula and invariance of the Haar integral under the compact rotation imply

\begin{equation}
\mathbb E\exp(tD_f\mathcal A)
\leq\exp\left[\tfrac12 b t^2\sum_e(f_y-f_x)^2\right].
\label{eq:pre:2-3}
\end{equation}

Here \(E(f)=\sum_e(f_y-f_x)^2\) and \(q^{tf}_x=e^{atf_x}q_x\).
Use \(\mathcal A(q^{-tf})\leq\mathcal A(q)-tD_f\mathcal A(q)+\frac12bt^2E(f)\), exponentiate, and integrate. The same Taylor estimate bounds the \(r\)-th moment of the likelihood ratio of a unit rotation by \(\exp(C_r bE(f))\). It is therefore bounded when \(bE(f)\) is bounded. These arguments use \(|w_et_e|\leq b\), so remain valid for signed and twisted interactions.

\subsubsection{Extracting a transverse current on a buffered patch}\label{pre:part-2-2}

Take an interior patch of lattice side \(d\), with a proportional buffer. For a scalar test \(h\) on the inner patch, let \(H\) be its \(C^1\) norm in patch coordinates. Choose a compact scalar potential \(\vartheta\), affine with slope \(1/d\) in direction \(\mu\) on all edges meeting the support of \(h\), and with bounded Dirichlet energy. Choose \(a\) perpendicular to the current color to be extracted. The compact score in the other perpendicular color has every fixed moment \(O_p(\sqrt b)\) by \eqref{eq:pre:2-3}.

Assume \(H>0\), since otherwise the test is zero. Rotate this score by the two opposite profiles

\[
 f=\tfrac12\arcsin(h/N),\qquad N=C\sqrt{1+b}\,H.
\]

The rotations have uniformly bounded likelihood norms. The difference of the rotated scores is the transverse current multiplied by \(2\sin(f_x+f_y)\). Multiplication by \(N\) extracts \(h\) at the midpoint up to \(CH/d\). It is important that only this first-order error is needed: after the score normalization \(d^{-1}\), its deterministic total over \(O(d^2)\) edges is \(O(bH)\). The moment cost of the two scores multiplied by \(N\) is \(O_p((1+b)H)\). Hence

\begin{equation}
\left\|d^{-1}\sum_{e\parallel\mu}w_e(s_e)_k h_e\right\|_p
\leq C_p(1+b)H.
\label{eq:pre:2-4}
\end{equation}

Bounded \(d\) can be treated deterministically.

\subsubsection{Whitney summation}\label{pre:part-2-3}

A width \(O(1)\) boundary layer contributes \(O(b\|F\|_\infty)\). Decompose the remaining square into buffered patches with \(d\) comparable to distance to the boundary. At one dyadic scale there are \(O(S/d)\) such patches; their normalized tests have \(C^1\) norms \(O(\|F\|_{C^1})\). Applying \eqref{eq:pre:2-4}, multiplying by \(d/S\), and summing gives \(O((1+b)\|F\|_{C^1})\) per scale. There are \(O(\log(2+S))\) scales. This proves \eqref{eq:pre:2-2}.
\end{proof}


\subsection{Conditioning into cells}\label{pre:cells}

The symbol \(L\) in this preliminary argument is an auxiliary integer
which may change during the stiffness iteration. It is independent of
the fixed block size chosen for the exact renormalization map in
Section~\ref{free:section}.

The current-moment estimate permits replacing slowly varying fields by
their values at cell anchors. We next express the large-square stiffness
through the conditioned cell laws; the variance of their conditional
currents is the gain to exploit.

Tile \(Q(S)\) by \(L^2\) cells of side \(R=S/L\). Condition on their complete boundary rings \(\omega\). Cell interiors are independent and the ring density has action
\(S_g(\omega)=-\sum_c\log Z_c(\omega)\). Set

\[
j(F)=\mathbb E[X_Q(F)\mid\omega]
=L^{-1}\sum_c\mathbb E_cX_c(F).
\]

Let \(z_c\) be the lower-left cell anchor and put

\[
N_c(F)=\sup_{c^+}(|F|+L^{-1}|\nabla F|),\qquad
D_c(F)=L^{-1}N_c(\nabla F).
\]

The fixed enlargement \(c^+\) is used only for smooth test-field estimates. Write \(P(B)\) for a fixed polynomial whose degree and coefficients may depend on fixed moment orders or prescribed exponents, but not on \(R,L,S\). Equation \eqref{eq:pre:2-2} and \eqref{eq:pre:2-1} give

\begin{equation}
|\mathbb E_cX_c(F)|\leq P(B)N_c(F),\quad
|\mathcal H_c(F,G)|\leq P(B)N_c(F)N_c(G),
\label{eq:pre:3-1}
\end{equation}

\begin{equation}
|\mathcal H_c(F,G)-\overline{F_c}H_cG_c|
\leq P(B)[D_c(F)N_c(G)+N_c(F)D_c(G)].
\label{eq:pre:3-2}
\end{equation}

For the last estimate subtract the constants \(F_c,G_c\) in the bilinear form. Their remainders have cell-coordinate \(C^1\) norm bounded by the displayed \(D_c\). One can instead use a direct Lipschitz bound on the cell where appropriate.

For a scalar ring rotation \(u\), let \(H_c^u\) be the resulting constant-field stiffness and \(O_c(u)=\operatorname{Ad}(e^{au(z_c)})\). Then

\begin{equation}
\|O_c(u)^TH_c^uO_c(u)-H_c\|\leq P(B)D_c(u),
\label{eq:pre:3-3}
\end{equation}

and its derivative at zero satisfies the same bound. To prove this, remove the constant rotation \(u(z_c)\) exactly by covariance. Along the remaining path also rotate the interior integration variables. The logarithmic density derivative is the negative of the centered cell current with test \(aR(u_y-u_x)\), whose fixed moments are bounded by \(P(B)D_c(u)\). The residual profile need not vanish on the cell boundary; the current-moment bound applies to this test without that restriction. Derivatives of a constant-field current have the bracket current with residual profile of size \(D_c(u)\), plus a deterministic contact of that size. Differentiate the contact and covariance in \eqref{eq:pre:1-1} and apply H\"older with the already established current moments. Restarting this calculation at each point of the path gives the bound for arbitrary amplitude \(u\); there is no exponential dependence on its constant part.

Rotating all rings while fixing the outside boundary gives the exact Hessian identity

\begin{equation}
D_u^2S_g=L^{-2}\sum_c\mathcal H_c(a\,d_Su,a\,d_Su).
\label{eq:pre:3-4}
\end{equation}


\subsection{Gaussian augmentation and commuting derivatives}\label{pre:gaussian}

The cell stiffness matrices need not be positive. We add auxiliary
Gaussian noise so that the desired upper bound on $H_Q$ becomes a lower
bound on the variance of an augmented current. The identity below makes
this reduction exact.

Suppose uniformly throughout the pinned cell class
\(H_c\leq hI\), where \(B^{-D}\leq h\leq B\). Choose

\begin{equation}
h<v\leq2B,\qquad c_0\geq B^{-D'},\qquad
C_c=vI-H_c-c_0P_a\geq B^{-D'}I,
\label{eq:pre:4-1}
\end{equation}

where \(P_a\) is projection onto the two spatial components of color \(a\). Let \(n_c\) be independent standard six-dimensional Gaussian vectors, and \(n_{0,c}\) independent two-dimensional Gaussian vectors of covariance \(c_0I\), all independent of the true ring data. Define

\begin{equation}
J(F)=j(F)+L^{-1}\sum_c(\sqrt{C_c}n_c+a n_{0,c})\cdot F_c.
\label{eq:pre:4-2}
\end{equation}

For every constant real matrix \(A\), total variance gives exactly

\begin{equation}
H_Q(A,A)=v|A|^2-\operatorname{Var}J(A).
\label{eq:pre:4-3}
\end{equation}

Indeed \(H_Q=L^{-2}\sum_c\mathbb EH_c-\operatorname{Var}j\), and the conditional noise covariance adds \(vI-L^{-2}\sum_cH_c\).

Choose a real trigonometric cutoff \(\psi\), zero on the boundary, with band \(O(m)\). Let \(f_i=\psi\varphi_i\), where \(\varphi_i\) are the real orthonormal Fourier modes of frequency \(|k|\leq4K\), including the constant. Use \(K+m\ll L\) and

\[
G_0=L^{-2}\sum_c\nabla f(z_c)\nabla f(z_c)^T,
\quad
p=(c_0G_0)^{-1}L^{-1}\sum_c\nabla f(z_c)n_{0,c},
\quad \theta=f\cdot p.
\]

Quadrature is exact at this band, so \(G_0\) is the continuum gradient Gram matrix. It is positive: a linear combination of the \(f_i\) with zero gradient is a constant, zero on the boundary, and hence is zero; division by \(\psi\) on its open nonzero set then annihilates the Fourier polynomial. Thus \(p\) is Gaussian with covariance \((c_0G_0)^{-1}\).

Put

\[
\omega^0=e^{-a\theta}\omega,\qquad
n_{0,c}^0=n_{0,c}-c_0\nabla\theta(z_c)/L,
\]

\[
n_c^0=O_c(\theta)^Tn_c-\sqrt{C_c^0}\,a\nabla\theta(z_c)/L,
\]

where \(C_c^0\) uses \(\omega^0\). At fixed base data \((\omega^0,n_0^0,n^0)\), the density of \(p\) is proportional to \(e^{-V(p)}\), with

\begin{equation}
V(p)=\tfrac12c_0p^TG_0p+S_g(e^{a\theta}\omega^0)
+\tfrac12\sum_c|n_c^0+\sqrt{C_c^0}\,a\nabla\theta(z_c)/L|^2.
\label{eq:pre:4-4}
\end{equation}

This is an exact change of variables. First split the Gaussian \(n_0\) into its linear projection and orthogonal residual. At fixed \(p\), the ring rotation preserves product Haar measure. At fixed \(p,\omega^0\), the \(n_c\) transformations are orthogonal maps followed by translations. These successive maps have constant Jacobian; dependencies of later maps on earlier variables give a triangular Jacobian, not an extra determinant.

For fixed scalar directions \(u_i\) in the potential span, let \(\partial_i\) mean addition of \(tu_i\) to \(\theta\) with all base data fixed. They commute. Conditional integration by parts gives, for \(\partial_i^*=-\partial_i+\partial_iV\),

\begin{equation}
\mathbb E\left(\sum_i\partial_i^*Y_i\right)^2
=\mathbb E\sum_{ij}(\partial_jY_i)(\partial_iY_j)
+\mathbb E\sum_{ij}Y_i(\partial_i\partial_jV)Y_j.
\label{eq:pre:4-5}
\end{equation}

The identity first holds conditionally and then in the full law. The conditional density has Gaussian decay: its first term is strictly positive quadratic and \(S_g\) is bounded for each fixed finite system. The functions involved and their derivatives have at most polynomial growth in Gaussian variables, since the matrices remain in a fixed positive cone for that system. Thus the integration by parts has no missing boundary term.


\subsection{Three uniform band estimates}\label{pre:bands}

The integration-by-parts identity in the preceding subsection contains
a Hessian term and a crossed derivative term. We next control their
ingredients uniformly in the number of modes, while retaining the exact
contact cancellation in the Ward identity.

Fix \(D,D'\) and any requested power \(N\). Take \(K/L\leq B^{-J}\), with \(J\) sufficiently large depending only on these fixed numbers and on the fixed test amplitude bounds. Retain \(S\geq L\), \eqref{eq:pre:2-1}, and \(Bm\leq K\). The constants denoted \(P(B)\) below do not depend on the subsequently increased \(J\).

\subsubsection{Sampling, gradients, and the good event}\label{pre:part-5-1}

For any trigonometric field of band \(O(K)\), with \(K\ll L\), we have

\begin{equation}
\|N(F)\|_\square\leq C\|F\|_2,\qquad
\|D(F)\|_\square\leq C(K/L)\|F\|_2,
\quad
\|a\|_\square^2=L^{-2}\sum_c|a_c|^2.
\label{eq:pre:5-1}
\end{equation}

Here is a direct sampling proof. Choose a smooth frequency cutoff equal to
one on the allowed band and supported in a ball of radius \(cL\), with fixed
\(c\). Its periodic convolution kernel \(K_L\) satisfies, for any fixed
\(N_0>2\),
\[
 |\partial^r K_L(z)|\leq C_{r,N_0}L^{2+r}
 (1+L\operatorname{dist}(z,\mathbb Z^2))^{-N_0}.
\]
This follows by writing the kernel as the periodization of the rescaled
inverse Fourier transform of the smooth cutoff and integrating that
transform by parts. Since \(F=K_L*F\), weighted Cauchy gives
\[
 \sup_{z\in c^+}|F(z)|^2\leq
 CL^2\int (1+L\operatorname{dist}(z_c-y,\mathbb Z^2))^{-N_0}|F(y)|^2\,dy.
\]
Sum over cells and divide by \(L^2\); the remaining lattice sum is bounded.
Apply the same argument to derivatives and use Parseval,
\(\|\partial^rF\|_2\leq(CK)^r\|F\|_2\). This proves
\eqref{eq:pre:5-1}, including its constant independent of the number of cells.
The centered-difference Fourier multiplier also gives

\begin{equation}
\|d_Su-\nabla u\|_2\leq C(K/S)^2\|\nabla u\|_2.
\label{eq:pre:5-2}
\end{equation}

Although the matrix \(G_0\) need not be well conditioned in the chosen basis, that is immaterial. The covariance of \(\nabla\theta\) is \(c_0^{-1}\) times the orthogonal projection onto the gradient subspace. Evaluation of a band \(O(K)\) gradient field has norm at most \(CK\), and evaluation after one more derivative at most \(CK^2\). Therefore

\begin{equation}
\rho=L^{-1}\|\nabla\theta\|_\infty+L^{-2}\|\nabla^2\theta\|_\infty
\leq B^{-n}
\label{eq:pre:5-3}
\end{equation}

outside an event of probability at most \(e^{-B^3}\), for any prescribed fixed \(n\) if \(J\) is large enough. For the supremum bound one can use only the elementary Fourier estimate
\(\|\nabla F\|_\infty\leq CK^2\|F\|_2\leq CK^2\|F\|_\infty\)
on this unit-volume torus. A grid of spacing \(cK^{-2}\) therefore captures
at least half the supremum. It has \(O(K^4)\leq e^{CB}\) points by
\eqref{eq:pre:2-1}. At each point the gradient standard deviation is at most
\(CK/\sqrt{c_0}\), and that of one more derivative at most
\(CK^2/\sqrt{c_0}\). After division by \(L\) and \(L^2\), respectively,
the Gaussian tail exponent at the threshold in \eqref{eq:pre:5-3} is an
arbitrarily high fixed power of \(B\) when \(J\) is increased. This proves
the claimed probability after the finite union bound.

\subsubsection{Hessian estimate}\label{pre:part-5-2}

From \eqref{eq:pre:4-4} and \eqref{eq:pre:3-4}, the exact conditional Hessian is the sum of the cell stiffness form and the two quadratic noise forms. Replace \(d_Su\) by \(\nabla u(z_c)\) using \eqref{eq:pre:3-2}, \eqref{eq:pre:5-1}, and \eqref{eq:pre:5-2}. Conjugate each true cell stiffness back to its base value using \eqref{eq:pre:3-3}. The color \(a\) is fixed by \(O_c\). The total error on \eqref{eq:pre:5-3} is at most

\[
P(B)[K/L+(K/S)^2+\rho]\,\|\nabla u\|_2^2.
\]

The remaining matrix sum is exactly \(H_c^0+C_c^0+c_0P_a=vI\). Consequently, for all directions simultaneously,

\begin{equation}
\partial_u^2V\leq(v+B^{-N})\|\nabla u\|_2^2
\label{eq:pre:5-4}
\end{equation}

on the good event. Everywhere the same upper bound holds with \(P(B)\) in place of \(v+B^{-N}\), by \eqref{eq:pre:3-1}, \eqref{eq:pre:4-1}, and \eqref{eq:pre:5-1}. The crude bound

\begin{equation}
\|J(F)\|_p\leq C_pP(B)L\|F\|_2
\label{eq:pre:5-5}
\end{equation}

follows from \eqref{eq:pre:3-1}, Cauchy over cells, and conditional Gaussian moments. Thus an exceptional event of probability \(e^{-B^3}\) contributes an arbitrary inverse power of \(B\) to every later finite moment expression, even if that expression has a fixed polynomial number of modes and a fixed polynomial factor in \(L\): \(\log L\leq16B\).

\subsubsection{Nonabelian derivative estimate}\label{pre:part-5-3}

Let at most \(CK^2\) real directions \(u_i\) satisfy

\begin{equation}
\|\nabla\sum_i d_i u_i\|_2\leq C|d|,\qquad
N_c(\nabla u_i)\leq C.
\label{eq:pre:5-6}
\end{equation}

For real deterministic \(F\) of band \(O(K)\), perpendicular in color to \(a\), write \(F'=-a\times F\). Then

\begin{equation}
\mathbb E\sum_i|\partial_iJ(F)-2J(u_iF')|^2
\leq B^{-N}\|F\|_2^2.
\label{eq:pre:5-7}
\end{equation}

Here is the termwise proof. For \(u=\sum d_iu_i\), differentiation of the physical conditional mean gives exactly

\begin{equation}
\partial_u j(F)=2j(uF')+j(r_uF')
+L^{-2}\sum_c\mathcal H_c(a\,d_Su,F),
\label{eq:pre:5-8}
\end{equation}

where \((r_u)_\mu(z)=u(z+e_\mu/(2S))+u(z-e_\mu/(2S))-2u(z)\). Taylor or Fourier multipliers give
\(\|r_u\|_2\leq CK S^{-2}\|\nabla u\|_2\).
Using cell products in \eqref{eq:pre:3-1} and \eqref{eq:pre:5-1},

\begin{equation}
|j(r_uF')|\leq P(B)LK S^{-2}\|\nabla u\|_2\|F\|_2.
\label{eq:pre:5-9}
\end{equation}

The anchor-conjugation part of differentiating \(\sqrt{C_c}\) combines with the derivative of \(n_c=O_c(\theta)(n_c^0+\sqrt{C_c^0}a\nabla\theta/L)\) to give exactly the noise part of \(2J(uF')\). The \(n_0\) noise is in color \(a\) and pairs to zero. The remaining deterministic noise drift is

\begin{equation}
L^{-2}\sum_c F_c\cdot\sqrt{C_c}\,O_c(\theta)\sqrt{C_c^0}\,a\nabla u(z_c).
\label{eq:pre:5-10}
\end{equation}

On the good event, \eqref{eq:pre:3-3} and the positive lower bound in \eqref{eq:pre:4-1} imply
\(\sqrt{C_c}O_c\sqrt{C_c^0}a=C_ca+O(P(B)\rho)\).
Together with the anchored version of the last term in \eqref{eq:pre:5-8}, this leaves \(F_c\cdot(H_c+C_c)a\nabla u\), which is zero because \(H_c+C_c=vI-c_0P_a\) and \(F\perp a\). The error is at most

\begin{equation}
P(B)[K/L+(K/S)^2+\rho]\|\nabla u\|_2\|F\|_2.
\label{eq:pre:5-11}
\end{equation}

Both \eqref{eq:pre:5-9} and \eqref{eq:pre:5-11} are bounds on a \emph{real scalar linear functional of \(d\)}, holding pointwise uniformly for \(|d|\leq1\). Its squared norm is precisely the sum of the squared errors in the individual directions. One must not multiply these bounds by the number of directions. Complex test fields follow by applying the estimate separately to their real and imaginary parts, at the cost of a fixed factor.

There remains the non-anchor square-root derivative \(\Gamma_c(u_i)\), with
\(\|\Gamma_c(u_i)\|\leq P(B)L^{-1}N_c(\nabla u_i)\).
At the true data the corresponding error is
\(L^{-1}\sum_cF_c\cdot\Gamma_c(u_i)n_c\).
The noises \(n_c\) are independent of \(\omega,p\); \(\Gamma_c(u_i)\) depends on \(\omega\) and the deterministic direction, not on these noises. Hence

\begin{equation}
\sum_i\mathbb E\left|L^{-1}\sum_cF_c\cdot\Gamma_c(u_i)n_c\right|^2
\leq P(B)\frac{K^2}{L^2}\|F\|_2^2.
\label{eq:pre:5-12}
\end{equation}

This last estimate does count the \(O(K^2)\) directions; its smallness comes from \(K/L\). On the exceptional angular event use the everywhere polynomial drift bound, and then its probability. Equations \eqref{eq:pre:5-9}--\eqref{eq:pre:5-12} prove \eqref{eq:pre:5-7} after choosing \(J\) sufficiently large. They explain exactly where a dimension loss is and is not allowed.

\subsubsection{Same-frequency Ward bound}\label{pre:part-5-4}

For the uncentered defect

\[
\mathcal D(F,G)=vL^{-2}\sum_c\overline{F_c}\cdot G_c
-\mathbb E\overline{J(F)}J(G),
\]

total covariance splitting and \eqref{eq:pre:3-2} compare it to the raw defect from \eqref{eq:pre:1-5}, with error

\begin{equation}
P(B)(K/L)\|F\|_2\|G\|_2.
\label{eq:pre:5-13}
\end{equation}

Let \(\eta=r\xi(z)e^{ik\cdot z}/(i|k|)\), with \(\eta\) zero on the boundary, \(|r|\leq1\), \(m\leq|k|\leq CK\), and \(\xi\) having band \(O(m)\), bounded supremum, and fixed polynomial \(C^3\) bounds. Let \(G\) be an exact gradient of a potential of the same type and frequency, or a bounded constant matrix times \(\xi_2e^{ik\cdot z}\). Require the amplitude's first derivative divided by \(|k|\) bounded, as is the case in all uses below. Apply \eqref{eq:pre:1-5} with \(\eta\) as defined; the identity already conjugates its first argument. Its contact cancels the raw defect. In the bracket term the two Fourier phases cancel even at the centered endpoints; the remaining current test has \(C^1\) norm \(P(B)/|k|\). For a gradient \(G\), its centered amplitude differences and the factors \(S\sin(k_\mu/(2S))\) have the same bound after division by \(|k|\). Thus \eqref{eq:pre:2-2} and \eqref{eq:pre:5-13} imply

\begin{equation}
|\mathcal D(d_S\eta,G)|\leq P(B)/|k|+B^{-N}.
\label{eq:pre:5-14}
\end{equation}

The raw Ward identity must be applied before replacing the gradient by its leading Fourier symbol. Otherwise the exact contact cancellation is lost.


\subsection{A weighted current estimate}\label{pre:weighted}

The band estimates control differentiated currents. We now sum them with
frequency weights whose convolutions remain summable. This is the
auxiliary estimate needed to bound the crossed derivative term in the
variance gain; it is not yet the stiffness improvement itself.

Fix an integer \(N_*>100(1+D+D')\) and apply the band estimates with accuracy
\(B^{-10N_*}\). Increase the fixed exponent \(J\) whenever required below.
The error estimates in the next two subsections are then bounded by
\(C B^{-N_*}\). To track this convention, all remaining weighted frequency
sums are bounded by fixed powers of \(T,T_w\leq CB\); all cutoff derivative
and cell-moment constants are fixed powers of \(B\). The only crude factors
in \(L\) occur with an \(S^{-2}\) Taylor remainder or an event of probability
\(e^{-B^3}\). Since \(S\geq L\) and \(\log L\leq16B\), those errors are
also at most \(B^{-N_*}\). No unweighted count of the main frequency band
is used in this convention.

Choose a sufficiently large fixed exponent \(J\) and set

\[
M=B^J,\quad K=L/B^J,\quad L\geq B^{10J},
\quad
\psi(z)=\prod_{\mu=1}^2[1-\cos^{2m}(\pi z_\mu)],
\quad m=\lceil B^4\rceil.
\]

All fixed powers of \(\psi\) have Fourier \(\ell^1\) norm bounded independently of \(B\), band \(O(m)\), and fixed derivative bounds polynomial in \(B\). Put \(\chi=\psi^2\); then

\begin{equation}
I_\chi=\int_{[0,1]^2}\chi^2\geq1-CB^{-2}.
\label{eq:pre:6-1}
\end{equation}

This follows from the Gaussian decay of \(\cos^{2m}(\pi z)\) away from the seams and an exceptional width \(O(m^{-1/2})\).

In a long step use \(K_1=L/B^{4J}\), \(K_2=L/B^{3J}\). In a short step use \(K_1=8M\), \(K_2=256M\). The main reciprocal lattice is \(2\pi\mathbb Z^2\), and

\begin{equation}
T=\sum_{M\leq|k|\leq K_1}|k|^{-2}
=\frac1{2\pi}\log(K_1/M)+O(M^{-1}),\qquad c\leq T\leq CB.
\label{eq:pre:6-2}
\end{equation}

The coefficient is \((2\pi)^{-2}\times2\pi=1/(2\pi)\). The sum here includes both signs of the Fourier frequencies, equivalently both normalized real modes for each pair.

\subsubsection{The weight and its convolution}\label{pre:part-6-1}

For dyadic \(H\) from \(M\) up to \(K_2\), let

\[
p_H(k)=H^{-2}e^{-\sqrt{|k|/H}},\quad
w(k)=\sum_Hp_H(k),\quad T_w=\sum_kw(k).
\]

Then

\begin{equation}
T_w\asymp1+\log(K_2/M),\quad
w*w\leq CT_ww,\quad
(m+|k|)^2w(k)\leq C,
\label{eq:pre:6-3}
\end{equation}

and \(w(k)\geq c(M+|k|)^{-2}\) for \(|k|\leq K_2/2\). Shifts of size \(O(m)\) change \(w\) by at most a fixed factor.

For completeness, if \(H\leq H'/4\),

\[
\sqrt{|l|/H}+\sqrt{|k-l|/H'}
\geq\tfrac12\sqrt{|l|/H}+\sqrt{|k|/H'},
\]

so summing \(l\) proves \(p_H*p_{H'}\leq Cp_{H'}\). If \(H'/4\leq H\leq H'\), split into \(|l|\leq|k-l|\) and its complement and use

\[
\sqrt{|l|}+\sqrt{|k-l|}\geq\sqrt{|k|}+c\sqrt{\min(|l|,|k-l|)}.
\]

The retained exponential is summable on the comparable scale. Summing over the dyadic pair of scales proves the convolution bound. The other claims follow directly by dyadic summation.

The tail \(|k|>K-O(m)\) is negligible to every fixed inverse power of \(B\), even against fixed polynomial factors in \(L,B,|k|/L\). In the long case \(\sqrt{K/K_2}=B^J\), whereas \(\log L\leq16B\); the short case has an even larger separation.

\subsubsection{Closing a two-color estimate}\label{pre:part-6-2}

Fix perpendicular colors \(r,r'\), with \(r'=-a\times r\). For both real modes in each frequency pair with \(2M\leq|k_p|\leq4M\), put

\[
b_p=|k_p|^{-2},\qquad
G_p^r=r\,d_S(\psi\widetilde\varphi_p/|k_p|),
\]

where the quadrature modes are \(\sqrt2\sin,-\sqrt2\cos\) for \(\sqrt2\cos,\sqrt2\sin\), respectively. The Ward bound gives

\begin{equation}
\sum_pb_p\mathbb E J(G_p^r)^2\leq Cv.
\label{eq:pre:6-4}
\end{equation}

Indeed the summed contact is \(O(v)\), the annular sum of \(b_p\) is \(O(1)\), and the error \(P(B)/M+B^{-N}\) can be made \(o(v)\).

Take directions \(u_i=\sqrt{w(k_i)}\psi\varphi_i\), \(|k_i|\leq K\), and set

\[
Y_{pi}=J(G_p^ru_i),\qquad
U^r=\sum_pb_p\sum_i\mathbb E Y_{pi}^2.
\]

The directions obey \eqref{eq:pre:5-6} by \eqref{eq:pre:6-3}, Parseval,
and the bounded multiplier norms of \(\psi\). Use the adjoints for the
base-data-fixed conditional law in \eqref{eq:pre:4-5}, and define
\[
 D_p=\sum_i\partial_i^*Y_{pi},\qquad F_p=-G_p^{r'}.
\]
The sign is fixed by \((r')'=-r\): the prime field of \(F_p\) is
\(G_p^r\). Conditional integration by parts and
\eqref{eq:pre:5-7} therefore give
\[
 \sum_pb_p\mathbb E[J(F_p)D_p]
 =\sum_{pi}b_p\mathbb E[(\partial_iJ(F_p))Y_{pi}]
 =2U^r+O(B^{-N_*})\sqrt{U^r}.
\]
Weighted Cauchy and \eqref{eq:pre:6-4} bound the square of this expression
by \(Cv\sum_pb_p\mathbb ED_p^2\). To estimate the latter sum, apply
\eqref{eq:pre:4-5} to each vector \(Y_p\). The Hessian contribution is
at most \(CvU^r+O(B^{-N_*})\). Bound the crossed derivative term by
\(\sum_{pij}b_p\mathbb E|\partial_jY_{pi}|^2\). A second use of
\eqref{eq:pre:5-7} bounds this Frobenius sum by

\begin{equation}
C\sum_{pij}b_p\mathbb E|J(G_p^{r'}u_iu_j)|^2+O(B^{-N_*})
\leq CT_wU^{r'}+O(B^{-N_*}).
\label{eq:pre:6-5}
\end{equation}

For the last inequality use complex Fourier modes in \(i,j\), apply \(w*w\leq CT_ww\), and remove the extra power of \(\psi\) by its bounded Fourier \(\ell^1\) norm and the shift bound for \(w\). Frequencies beyond the truncated set are negligible by the preceding tail estimate and the elementary extension of \eqref{eq:pre:5-5}, \(\|J(G_p^{r'}\psi e^{ikz})\|_2\leq P(B)L(1+|k|/L)\). The derivative-error sum costs \(B^{-N}\sum_i\|G_p^ru_i\|_2^2\leq CB^{-N}T_w\), not \(B^{-N}K^2\).

We have consequently proved
\[
 \sum_pb_p\mathbb ED_p^2
 \leq CvU^r+CT_wU^{r'}+O(B^{-N_*}).
\]
Interchanging the colors gives the same inequality with \(r,r'\)
reversed, since \(U^{-r}=U^r\).
Thus neither color estimate is assumed in proving the other.
Let \(U=\max(U^r,U^{r'})\). Unless \(U\leq CB^{-2N_*}\), the signal lower
bound gives
\(U^2\leq Cv(v+T_w)U+CvB^{-N_*}\).
Solving this quadratic and using \(v\geq B^{-D}\) gives

\begin{equation}
U^r+U^{r'}\leq Cv(v+T_w).
\label{eq:pre:6-6}
\end{equation}

\subsubsection{Extracting all spatial directions}\label{pre:part-6-3}

At fixed total complex frequency \(n=k_p+k_i\), \(|n|\leq4K_1\), take the common positive lower bound
\(w(n-k_p)\geq c(M+|n|)^{-2}\). Write
\[
 \psi_\mu^+(z)=\frac12\bigl[\psi(z+e_\mu/(2S))+
                         \psi(z-e_\mu/(2S))\bigr].
\]
The exact gradient factor in spatial direction \(\mu\) is

\[
\widehat k_{p\mu}\operatorname{sinc}(k_{p\mu}/(2S))\psi_\mu^+
+\frac{S\cos(k_{p\mu}/(2S))}{i|k_p|}
[\psi(z+e_\mu/(2S))-\psi(z-e_\mu/(2S))].
\]

Under \(k_p\mapsto-k_p\) its first part changes sign and its second part does not. Pairing the two squared currents cancels the interference term and permits discarding the second square. Separate coordinate reflections of \(k_p\) then cancel the mixed spatial terms. The annular sum of each diagonal coefficient is bounded below by a positive constant. Replacing \(\psi\psi_\mu^+\) by \(\psi^2\) has Fourier size \(P(B)S^{-2}\), so \eqref{eq:pre:5-5} makes its summed error negligible. It follows that

\begin{equation}
\boxed{\quad
W^r:=\sum_{|n|\leq4K_1}(M+|n|)^{-2}
\sum_{\mu=1}^2\mathbb E|J(r e_\mu\psi^2e^{in\cdot z})|^2
\leq Cv(v+T_w).\quad}
\label{eq:pre:6-7}
\end{equation}

This is the auxiliary weighted estimate. It has not been assumed in any preceding replacement.


\subsection{The variance gain}\label{pre:variance}

We now use the weighted estimate to obtain a lower bound on current
variance. Through \eqref{eq:pre:4-3}, this becomes the desired decrease
of the stiffness upper bound.

For a real constant \(3\times2\) matrix \(A\) of Hilbert--Schmidt norm one, choose a unit \(a\) perpendicular to both its columns. This is where the three-dimensional imaginary quaternion space is used. Put \(A'=-a\times A\), so \(|A'|=1\). On the main annulus set

\[
u_i=\chi\varphi_i/|k_i|,\quad
G_i^{A'}=d_S\bigl(\chi(A'\widehat k_i)\widetilde\varphi_i/|k_i|\bigr),\quad
Y_i=J(G_i^{A'})/|k_i|.
\]

The gradient Gram operator of the \(u_i\) is bounded by \((1+Cm/M)^2\). The Ward bound, exact quadrature, and
\(\sum|k|^{-2}\widehat k\widehat k^T=(T/2)I_2\) yield

\begin{equation}
\sum_i\mathbb E Y_i^2=vI_\chi T/2+O(B^{-N_*}),\qquad
\sum_i\mathbb E(\partial_iJ(A))Y_i=vI_\chi T+O(B^{-N_*}),
\label{eq:pre:7-1}
\end{equation}

\begin{equation}
\mathbb E Y^TV''Y\leq v^2I_\chi T/2+O(B^{-N_*}).
\label{eq:pre:7-2}
\end{equation}

The second formula in \eqref{eq:pre:7-1} uses the factor two in \eqref{eq:pre:5-7}. The leading-symbol replacement of \(G_i\) is made only inside the contact pairing after applying \eqref{eq:pre:5-14}, with error \(C[m/|k_i|+(|k_i|/S)^2]\). Summing with \(|k_i|^{-2}\) makes the total arbitrarily small. The same error convention, chosen stronger than all needed powers of \(v\), is used in this section.

\subsubsection{The crossed term}\label{pre:part-7-1}

The estimate needed is

\begin{equation}
\left|\sum_{ij}\mathbb E(\partial_jY_i)(\partial_iY_j)\right|
\leq C[vT^2+v(v+T_w)].
\label{eq:pre:7-3}
\end{equation}

By \eqref{eq:pre:5-7} the derivative array is, up to a negligible summed-square error,
\(-2J(G_i^A\chi\varphi_j)/(|k_i||k_j|)\). Its complex change of basis deserves a convention check. For a real matrix \(D\), transforming both indices by a unitary matrix \(U\) produces \(D'=UDU^T\). Then
\(\sum_{ij}D'_{ij}\overline{D'_{ji}}=\operatorname{Tr}(D^2)\), because \(U^T\overline U=I\). Thus the crossed pairing becomes the sesquilinear pairing of the two arrays at the \emph{same total frequency} \(n=k+l\); it is not a bilinear pairing without conjugation in the complex basis.

For main-annulus frequencies define

\[
H_0(n)=\sum_{k+l=n}|k|^{-2}|l|^{-2}.
\]

Elementary annular summation gives

\begin{equation}
H_0(n)\leq CT(M+|n|)^{-2}.
\label{eq:pre:7-4}
\end{equation}

For \(|n|<2M\), \(T\) can be replaced by a constant. For \(|n|\geq2M\), it can also be replaced by a constant if the summand is multiplied by
\(\min(1,\min(|k|,|l|)/|n|)\). To prove these statements split where \(|k|\leq|n|/2\), where \(|l|\leq|n|/2\), and the complement. In the first region the other denominator is \(O(|n|^{-2})\), and the annular sum of \(|k|^{-2}\) produces the logarithm; the inserted \(|k|/|n|\) removes it. The complementary comparable-size region and its \(|k|\geq2|n|\) tail are summable without a logarithm. For small \(n\), the cutoff \(|k|,|l|\geq M\) gives the bound directly.

Together with \eqref{eq:pre:6-7}, \eqref{eq:pre:7-4} bounds the summed square norms of fields \(\psi^2e^{inz}\) with bounded constant transverse-color coefficients by \(CTv(v+T_w)\). A Fourier multiplier whose coefficients have a common summable majorant of total \(C\), supported at shifts \(O(m)\), preserves this estimate. Cauchy with the absolute Fourier coefficients and the bounded shift cost of \((M+|n|)^{-2}\) proves this last claim.

It is therefore legitimate to replace \(G_k^A\chi e^{ilz}\) by

\begin{equation}
e^{inz}\chi^2A P_k,\qquad P_k=\widehat k\widehat k^T.
\label{eq:pre:7-5}
\end{equation}

Before and after replacement, factor out \(\psi^2e^{inz}\); the residual Fourier coefficients have common \(\ell^1\) majorants. The error majorant is \(C[m/M+(K_1/S)^2]\), from the exact centered-gradient formula. Hence all these replacements are negligible by \eqref{eq:pre:6-7}, \eqref{eq:pre:7-4}, and Cauchy. The small-\(n\) part already obeys \eqref{eq:pre:7-3}.

For \(|n|\geq2M\), split the spatial index in \eqref{eq:pre:7-5} into directions parallel and perpendicular to \(n\). Replace the longitudinal component by

\begin{equation}
d_S\left(\chi^2 A P_k\widehat n\,e^{inz}/(i|n|)\right).
\label{eq:pre:7-6}
\end{equation}

This replacement also has a small summed-square error, but two sorts of error must be distinguished. The sinc error and the first derivative of \(\chi^2=\psi^4\) are divisible by \(\psi^2\) and have common Fourier majorants of total \(C[(K_1/S)^2+m/M]\), so \eqref{eq:pre:6-7} applies. The symmetric-shift error and the centered Taylor remainder have Fourier size \(P(B)/S^2\) and need not be divisible by \(\psi^2\); use \eqref{eq:pre:5-5} instead. Its factor \(L/S^2\leq1/L\) makes their sum negligible. This avoids an unjustified division by a cutoff that vanishes at the boundary.

If at least one paired field is now the exact gradient \eqref{eq:pre:7-6}, \eqref{eq:pre:5-14} bounds its uncentered current pairing with the other same-frequency field by \(Cv\). Summing \(|k|^{-2}|l|^{-2}\) yields \(CvT^2\).

For two transverse spatial components, the product of coefficient norms is at most

\begin{equation}
|\sin\angle(k,n)|\,|\sin\angle(l,n)|
\leq\min\left(1,\frac{\min(|k|,|l|)}{|n|}\right).
\label{eq:pre:7-7}
\end{equation}

For example, if \(|k|\leq|l|\), then \(|\sin\angle(l,n)|=|k\times l|/(|l||n|)\leq|k|/|n|\); bound the other sine by one. Bound the pairing of the constant-coefficient fields by this product times the sum of the four coordinate current second moments. The improved version of \eqref{eq:pre:7-4} and \eqref{eq:pre:6-7} then give \(Cv(v+T_w)\). This proves \eqref{eq:pre:7-3}.

\subsubsection{Variance gain}\label{pre:part-7-2}

Set \(\mathcal S=\sum_i\partial_i^*Y_i\), using the directions and
vector field defined above \eqref{eq:pre:7-1}. Conditional
integration by parts gives \(\mathbb E[\mathcal S\mid\text{base data}]=0\)
and
\[
 \operatorname{Cov}(J(A),\mathcal S)
 =\sum_i\mathbb E[(\partial_iJ(A))Y_i]
 =vI_\chi T+O(B^{-N_*}).
\]
Equation \eqref{eq:pre:4-5} expresses \(\mathbb E\mathcal S^2\) as the
Hessian contribution plus the crossed derivative term. Their bounds
\eqref{eq:pre:7-2}--\eqref{eq:pre:7-3}, followed by
\(|\operatorname{Cov}(J(A),\mathcal S)|^2
\leq\operatorname{Var}J(A)\mathbb E\mathcal S^2\), give

\begin{equation}
\operatorname{Var}J(A)
\geq
\frac{(vI_\chi T+O(B^{-N_*}))^2}
{v^2I_\chi T/2+C[vT^2+v(v+T_w)]+O(B^{-N_*})}.
\label{eq:pre:7-8}
\end{equation}

In the long regime \(v\geq h\geq1\), \(\log L\leq h\), use \((1+x)^{-1}\geq1-x\), \eqref{eq:pre:6-1}, and \eqref{eq:pre:6-2}:

\begin{equation}
\operatorname{Var}J(A)
\geq2I_\chi T-C[1+T^2/v+T_w/v]
\geq\frac{\log L}{\pi}-C_J\left[\log B+\frac{(\log L)^2}{v}\right].
\label{eq:pre:7-9}
\end{equation}

In the short regime \(T,T_w\) are bounded above and below by absolute constants, and \eqref{eq:pre:7-8} gives

\begin{equation}
\operatorname{Var}J(A)\geq c_{\rm sh}\min(v,1).
\label{eq:pre:7-10}
\end{equation}

The positive constant \(c_{\rm sh}\) can be fixed before choosing the final large \(J\) and \(B\): all analytic leading constants in \eqref{eq:pre:6-7} and \eqref{eq:pre:7-3} are absolute, while those choices only suppress errors. This is necessary for the low-stiffness induction below.


\subsection{Iteration of the stiffness bound}\label{pre:iteration}

Fix any sufficiently large \(D\), then \(D'>D\), and fix \(J\) to make all preceding errors as small as required. Start with the uniform upper bound \(H_1\leq bI\). If \(b\leq B^{-D}\), no step is needed.

For \(h\geq A\log B\), use

\[
v=h+1,\quad c_0=1/2,\quad
L=\left\lceil e^{\sqrt{h\log B}}\right\rceil.
\]

Choose fixed \(A\) large compared with \(J^2\) and the constants in \eqref{eq:pre:7-9}. Then \(L\geq B^{10J}\), \(\log L\leq h\), \eqref{eq:pre:4-1} holds, and \eqref{eq:pre:4-3}, \eqref{eq:pre:7-9} decrease the stiffness upper bound by

\[
\delta h\geq(\log L)/\pi-C\log B
\geq c\sqrt{h\log B}.
\]

There are at most \(C\sqrt{B/\log B}\) such steps. Summing the relation \(\log L\leq\pi\delta h+C\log B\) gives a total logarithmic side length at most

\begin{equation}
\pi b+C_D\sqrt{B\log B}.
\label{eq:pre:8-1}
\end{equation}

Stop at \(h\leq C\log B\), using a weaker positive upper bound at the crossing if necessary. Now take a small fixed \(\delta<c_{\rm sh}/4\), and use

\[
v=h+\delta\min(h,1),\quad
c_0=\tfrac12\delta\min(h,1),\quad
L=\lceil B^{11J}\rceil.
\]

The matrix lower bounds in \eqref{eq:pre:4-1} hold with \(D'>D\), after increasing the fixed lower threshold for \(B\). Equations \eqref{eq:pre:4-3}, \eqref{eq:pre:7-10} reduce \(h\) by at least \(c\min(h,1)\). It takes \(O_D(\log B)\) steps to reach \(B^{-D}\), at total additional logarithmic length \(O_D((\log B)^2)\). This is absorbed in \eqref{eq:pre:8-1}.

At all stages \(\log(2+S)\leq16B\) holds once \(B\) exceeds a threshold depending only on the fixed choices: the accumulated bound is \(\pi b+o(B)\), and a proposed next step adds at most \(B\). This verifies the premise used in the moment bounds, rather than assuming it retrospectively.

We have proved, uniformly over the signed/twisted/pinned class, that there is a common integer side \(R_*\) with

\begin{equation}
\boxed{\ H_{R_*}\leq B^{-D}I,\qquad
\log R_*\leq\pi b+C_D\sqrt{B\log B}.\ }
\label{eq:pre:8-2}
\end{equation}


\subsection{Uniform control of local rotations}\label{pre:rotations}

Take coarse cells of integer side comparable to
\(s=R_*M_*\), where \(M_*=\lceil B^{D+C_0}\rceil\). Rectangles with sides in \([s,2s]\) allow arbitrary sufficiently large tori to be tiled. Profiles below are supported on a fixed number of coarse cells, vanish at the outer boundary of a fixed collar, are continuous across cells, and have uniformly bounded cellwise \(C^2\) extensions. Pins lie outside the collar.

For one fixed-axis profile \(f\), the effective coarse ring action \(F\) obeys

\begin{equation}
D_f^2F\leq CB^{-D}
\label{eq:pre:9-1}
\end{equation}

at every ring configuration. To prove it, subdivide each coarse cell into side-\(R_*\) squares and possible leftover strips of total area \(O(sR_*)\). Conditional variance subtraction bounds its Hessian above by the sum of the subcell Hessians. The gradient test on a small square has size \(O(M_*^{-1})\), with variation smaller by another \(M_*^{-1}\). Equations \eqref{eq:pre:8-2}, \eqref{eq:pre:3-2} bound each small-square cost by
\(CM_*^{-2}[B^{-D}+P(B)/M_*]\).
Summing gives \(CB^{-D}+P(B)/M_*\). On leftover strips discard the negative variance and use the contact bound: their total cost is \(CBR_*/s=CB/M_*\). Choose \(C_0\) large enough. The polynomial in this use of \eqref{eq:pre:3-2} is independent of \(D\), so the choice is not circular.

Haar integration by parts on the free ring spins gives

\begin{equation}
\mathbb E(D_fF)^2=\mathbb ED_f^2F\leq CB^{-D}.
\label{eq:pre:9-2}
\end{equation}

For the normalized square-root density \(\varphi\) relative to product Haar, this is \(\|D_f\varphi\|_2^2\leq CB^{-D}\). Pullbacks by deterministic ring rotations are unitary. Integrate the derivative along a fixed-axis rotation and telescope a bounded number of such factors. For every allowed product \(\phi\),

\begin{equation}
\mathbb E\left(e^{-[F(\phi\omega)-F(\omega)]/2}-1\right)^2\leq CB^{-D}.
\label{eq:pre:9-3}
\end{equation}

Thus a specified action displacement exceeds fixed \(\epsilon>0\) with probability \(O_\epsilon(B^{-D})\).

\subsubsection{Uniformity over profile parameters}\label{pre:part-9-1}

Let \(\phi_\lambda\) vary smoothly on a fixed-dimensional compact parameter rectangle, with fixed bounded spatial and parameter derivatives, and have a compact smooth homotopy to identity of this same class. Such a homotopy can be subdivided into finitely many small increments, each factored into three fixed-axis profiles by smooth local product coordinates on \(S^3\), so \eqref{eq:pre:9-3} holds uniformly at each specified parameter.

A compact microscopic Lie-algebra score with bounded Dirichlet energy is subGaussian with variance bound \(CB\), by \eqref{eq:pre:2-3} and H\"older over its three colors. The coarse score is its conditional expectation and has the same bound. The likelihood of any one of the present rotations is bounded in fixed \(L^p\) by \(e^{C_pB}\); apply the score and second-variation estimates for one factor and the likelihood cocycle and H\"older for finitely many factors. Consequently at rotated data

\[
\Pr(|\partial_\lambda F(\phi_\lambda\omega)|>t)
\leq\min(1,\exp[CB-ct^2/B]),
\]

In particular, choose a fixed \(A\) so that the exponent is at most
\(-ct^2/(2B)\) for \(t\geq AB\). The layer-cake formula then gives
\[
 \mathbb E|\partial_\lambda F(\phi_\lambda\omega)|^p
 \leq (AB)^p+p\int_{AB}^{\infty}t^{p-1}e^{-ct^2/(2B)}\,dt
 \leq C_p(1+B)^p.
\]
Thus the exponential likelihood bound produces a polynomial moment bound
after its Gaussian tail is taken into account.

Take \(p>2d\), where \(d\) is the fixed parameter dimension. Morrey's estimate bounds the parameter H\"older-\(1/2\) seminorm of the displacement by \(B^2\) with probability \(1-O(B^{-p})\). A mesh of spacing \(B^{-6}\) has \(O(B^{6d})\) points and then incurs a displacement error at most \(CB^{-1}\). Union bounding \eqref{eq:pre:9-3} over this mesh proves

\begin{equation}
\Pr\left(\sup_\lambda|F(\phi_\lambda\omega)-F(\omega)|>\epsilon\right)
\leq\zeta(B),\qquad \zeta(B)\longrightarrow0,
\label{eq:pre:9-4}
\end{equation}

provided \(D>10d+10\), for example. All estimates are uniform under conditioning outside the collar.


\subsection{Coarse orientations and decay of charged observables}\label{pre:orientations}

Independently adjoin Haar variables \(g_X\in S^3\) at free coarse vertices, equal to identity at pinned boundary vertices. On a directed coarse edge \(X\to Y\), interpolate by

\[
\Phi(g;t)=g_X\exp[\eta(t)\ell(g_X^{-1}g_Y)],
\]

where \(\eta\) is fixed smooth and constant near endpoints, and \(\ell\) is any fixed measurable logarithm of norm at most \(\pi\), equal to zero at identity. Write the retained spin ring as \(\omega=\Phi(g)\widetilde\omega\). At fixed \(g\) this is a Haar-preserving change of ring variables. Conditional on \(\widetilde\omega\), the \(g\)'s have density proportional to \(e^{-F(\Phi(g)\widetilde\omega)}\), a finite-range Gibbs density.

For every unpinned cell its potential is invariant under common left multiplication of its incident \(g\)'s: \((kg_X)^{-1}(kg_Y)=g_X^{-1}g_Y\), so \(\Phi(kg)=k\Phi(g)\), and simultaneous left rotation of the cell's rings and interior preserves its action even with the fixed right twists.

Call an interior vertex bad if, for some values of the neighboring \(g\)'s, varying its own \(g_X\) changes the incident-cell action by more than \(\epsilon\). Enlarge the test to permit all logarithms of the prescribed bounded size satisfying the endpoint equations, using the same data on unaffected edges. This enlargement makes the test a supremum over a compact parameter set and covers every branch choice of the measurable interpolation.

\subsubsection{Why the required profiles form finitely many smooth families}\label{pre:part-10-1}

Fix a star and its true redundant data. Two trial assignments differing only at its center produce relative ring rotations
\(\Phi(g^{(j)})\Phi(g^{\rm true})^{-1}\), \(j=1,2\). They agree along the star perimeter. Treat vertex quaternions and all bounded edge logarithms, including the true ones, as independent ambient parameters, imposing their endpoint equations only on the compact subset of valid assignments.

At any valid parameter value the finitely many relative edge arcs, together with identity, omit some point of \(S^3\): a finite union of smooth compact curves has zero three-dimensional volume. Their distance from that chosen point is positive and remains positive on a parameter neighborhood. Use stereographic coordinates on its complement, and ordinary local coordinates for the finitely many vertex parameters.

For neighboring ambient parameters not satisfying endpoint equations, correct each arc in stereographic coordinates near its endpoints so that it takes the prescribed relative vertex values. The corrections use fixed smooth cutoff functions and are zero on the valid subset. The original edge profiles are flat near vertices, so these corrected profiles are uniformly smooth within each coarse edge.

Extend to coarse cells by a Coons interpolation: add the two linear normal extensions of each pair of opposite side curves and subtract the bilinear interpolant of the four corner values. This matches all four sides. In the outer collar choose the additional edge and vertex data tending to identity, and choose this construction using only the common star-perimeter data. Thus the two extensions agree everywhere outside the star for valid assignments. The profiles are identity on and outside the outer collar boundary. Linear homotopy of their stereographic coordinates to the identity coordinate gives compact homotopies with the required uniform cellwise derivative bounds.

The compact valid parameter set has a finite subcover by such neighborhoods, which may be reduced to closed coordinate rectangles. All chart, extension, derivative, and homotopy constants are then finite and independent of the lattice side, coupling, spin configurations, and actual branch choices. Rectangular coarse-cell aspect ratios in \([1/2,2]\) add only compact parameters.

This justifies applying \eqref{eq:pre:9-4} to the entire bad-site test. In the original joint law, condition on the true \(g\)'s. The test is now a local event in the true ring variables \(\omega\), of probability at most \(\zeta\) even conditional on all rings outside a fixed collar. The two complete action displacements differ exactly by the incident-star displacement, since the extensions agree outside the star. A finite coloring of the collar-overlap graph selects at least \(|E|/C\) vertices with disjoint collars from any finite set \(E\). Iterated conditioning gives

\begin{equation}
\Pr(E\text{ is entirely bad})\leq\zeta^{|E|/C}.
\label{eq:pre:10-1}
\end{equation}

\subsubsection{Disagreement exploration with correct marginals}\label{pre:part-10-2}

The exploration follows the disagreement-percolation mechanism of
van den Berg and Maes~\cite{VanDenBergMaes1994}. Their finite-spin theorem
is not invoked for the present continuous spin space: the Haar
minorization and the adaptive marginal argument are proved below.

Fix \(\widetilde\omega\) and call the conditional orientation law \(\nu\). Its pushforward \(\nu_k\) under common left multiplication by \(k\) on all free vertices differs from \(\nu\) only in potentials near the pinned boundary. At each good interior vertex every one-site conditional density, under either law and every neighboring configuration, is at least \(e^{-\epsilon}\) times Haar. The same holds with only partial conditioning, by averaging the full conditional densities.

Seed a possible disagreement on a fixed-width boundary layer. Repeatedly reveal a vertex neighboring possible disagreements. At a good vertex couple the two appropriate conditional marginals by their common \(e^{-\epsilon}\)-Haar component; use a preassigned independent Bernoulli variable of parameter \(q=1-e^{-\epsilon}\) for the possible failure. At a bad vertex always permit failure. A successful vertex is assigned the same value. If the remaining unassigned components no longer touch a disagreement, their conditional distributions agree by the finite-range Markov form and can be sampled identically.

The adaptive order causes no marginal error. For a test \(H\) of one completed configuration, the original-law conditional expectation \(\nu[H\mid\text{that marginal's assigned values}]\) is a martingale relative to the full exploration history: for whichever vertex the history selects, sampling its value from that original-law conditional uses exactly the tower identity. This proof also covers the joint sampling of a remaining component at the stopping time.

A disagreement at a specified interior vertex requires an open simple path to the boundary seed, where all good vertices have failed their Bernoulli trials. The trials can be chosen independent of \(\widetilde\omega\). For a fixed path of \(r\) tested vertices, \eqref{eq:pre:10-1} gives

\begin{equation}
\mathbb E[q^{\#\mathrm{good}}]
\leq(q+\zeta^{1/C})^r.
\label{eq:pre:10-2}
\end{equation}

This follows by expanding \(\prod_v[q+(1-q)\mathbf1_{v\text{ bad}}]\) and applying \eqref{eq:pre:10-1} to every subset. Choose fixed \(\epsilon\) small and then \(B\) large enough that the right side beats the bounded-degree path count. The connection probability is then at most \(Ce^{-c\,\mathrm{distance}/s}\).

For a bounded original observable \(O\) supported on \(A\), whose average under a common left rotation is identically zero, integrate cell interiors at their transformed rings. The resulting \(O^*(g)\) depends on at most \(C|A|\) coarse vertices and still has identically zero common-left average. Comparing \(\nu\) with \(\nu_k\), using \eqref{eq:pre:10-2}, and averaging \(k\) proves

\begin{equation}
|\mathbb EO|\leq C\|O\|_\infty |A|e^{-c d_{\rm pin}/s}.
\label{eq:pre:10-3}
\end{equation}

For bounded Haar-centered one-site \(O(q_x)\), condition on the boundary of a coarse region within a third of the distance from \(x\) to another site \(y\). Apply \eqref{eq:pre:10-3} inside and multiply by any bounded one-site observable at \(y\). In unpinned sufficiently large tori,

\begin{equation}
|\mathbb E O(q_x)O'(q_y)|
\leq C\|O\|_\infty\|O'\|_\infty e^{-c d(x,y)/s}.
\label{eq:pre:10-4}
\end{equation}

Small distances use the trivial bound. Rectangular tilings and collars work whenever each periodic side is at least \(Cs\). Free subgraphs are obtained by padding with zero couplings. Right twists and signed couplings are still allowed in this step; no association is being used yet.


\subsection{Mixing for arbitrary local observables}\label{pre:neutral}

From now on use the untwisted zero-field ferromagnet, with \(0\leq b_e\leq b\). This restriction is essential. Neither the association argument nor its conclusion is being claimed for arbitrary pinned or frustrated systems.

\subsubsection{Association and the covariance inequality}\label{pre:part-11-1}

The covariance estimate below belongs to the association literature
\cite{Newman1980,BulinskiShabanovich1998}; the elementary proof is included.
For a finite associated collection of scalar random variables \(X_i\)
with finite second moments, and functions with individual Lipschitz
constants \(\ell_{F,i},\ell_{G,j}\),

\begin{equation}
|\operatorname{Cov}(F,G)|
\leq\sum_{ij}\ell_{F,i}\ell_{G,j}\operatorname{Cov}(X_i,X_j).
\label{eq:pre:11-1}
\end{equation}

Let \(U=\sum_i\ell_{F,i}X_i\), \(V=\sum_j\ell_{G,j}X_j\). The functions \(U\pm F,V\pm G\) are increasing. Adding the two matching-sign covariance inequalities gives the lower bound in \eqref{eq:pre:11-1}; adding the opposite-sign inequalities gives the upper bound.

The next doubled-variable proof is the plane-rotator form of Ginibre's
inequality~\cite{Ginibre1970}; it is not an application of a general
sphere-spin Ginibre inequality. For zero-field ferromagnetic XY, cosine-character means are nonnegative by Fourier expansion. The covariance of two cosine characters is nonnegative by the two-replica substitution \(\theta=\varphi+\psi\), \(\theta'=\varphi-\psi\). This map is an onto homomorphism of the product tori and pushes normalized Haar measure to normalized Haar measure. Each difference is
\(-2\sin(k\cdot\varphi)\sin(k\cdot\psi)\). Expanding the replicated weight
\(\exp(\sum_e2J_e\cos\varphi_e\cos\psi_e)\) therefore gives a sum of squares
of real integrals with nonnegative coefficients. Absolute convergence
follows from compactness and the exponential series. This proves the
required edge-energy covariance and monotonicity of two-point functions
in every coupling.

For zero-field ferromagnetic Ising, the random-cluster expansion
of Fortuin--Kasteleyn and its joint spin--bond form
\cite{FortuinKasteleyn1972,EdwardsSokal1988} with \(p_e=1-e^{-2J_e}\) gives a measure proportional to independent Bernoulli weights times \(2^{\#\mathrm{components}}\). Component count is supermodular because graph rank is submodular. Hence the edge law is associated, connection probabilities are nondecreasing in \(p_e\), and \(\langle\tau_i\tau_j\rangle=\Pr(i\leftrightarrow j)\). The spin law itself is associated by its elementary log lattice condition. These statements also hold with zero couplings by continuity.

The use of scalar amplitudes and conditional ferromagnets is related to
the amplitude-association argument of Campbell and Chayes for the $O(3)$
Wolff representation~\cite{CampbellChayes1998}; the two-circle $O(4)$
decomposition below is proved separately. The continuous angle marginals
below have a positive density satisfying the log lattice condition relative to product one-site measure: their mixed derivatives are nonnegative. This is the FKG lattice condition~\cite{FortuinKasteleynGinibre1971}.
For this continuous product space, association follows by induction in
the number of coordinates. Conditional densities at ordered values of one coordinate have an increasing density ratio; association on the remaining slice orders their increasing expectations, and one-dimensional Chebyshev's inequality completes the induction. Endpoint zero densities can be handled by restriction and a limit.

\subsubsection{A quantitative arbitrary-observable bound inside XY}\label{pre:part-11-2}

Split an XY spin as
\((\tau^1\cos\gamma,\tau^2\sin\gamma)\), \(\gamma\in[0,\pi/2]\). Under
the a priori circle measure the two signs are independent fair signs and
\(\gamma\) has density \(2/\pi\). Given the angles, the two sign species are
independent ferromagnetic Ising systems. The \(\gamma\) marginal is associated:
its density relative to the product angle measure is the product of the
two Ising partition functions. Differentiate its logarithm. Products of
first coupling derivatives have the same sign, edge covariances are
nonnegative, and mixed derivatives of a coupling on a single edge are
nonnegative.

Write \(c_{ij}=\langle s_i\cdot s_j\rangle\) for the full XY two-point function. For either sign species, and for the angle itself,

\begin{equation}
\mathbb E[\tau_i^h\tau_j^h]\leq Cc_{ij}^{1/3},\qquad
|\operatorname{Cov}(\gamma_i,\gamma_j)|\leq Cc_{ij}^{1/3}.
\label{eq:pre:11-2}
\end{equation}

To verify the second claim as well as the first, use an embedded Ising decomposition for any fixed reflection of the circle. A bounded one-site function odd under that reflection has the form \(\tau_i f_i\) after conditioning on its reflecting-component magnitude and perpendicular component. Its two-point is bounded by its sup norms times the averaged connection probability. By rotation invariance and the uniform one-site marginal,

\[
\tfrac12c_{ij}=\mathbb E[|x_i||x_j|\mathbf1_{i\leftrightarrow j}],
\]

where \(x_i\) is the reflecting coordinate. Since \(\Pr(|x_i|\leq\delta)\leq C\delta\),
\(\Pr(i\leftrightarrow j)\leq C\delta+C\delta^{-2}c_{ij}\).
Choose \(\delta\asymp c_{ij}^{1/3}\), with the endpoint cases by a limit or the trivial bound. Sign coordinates are odd under coordinate reflection; \(\gamma-\pi/4\) is odd under exchange of the two coordinates, itself a reflection. This proves \eqref{eq:pre:11-2}. The averaged conditional sign correlation is exactly the unconditional sign correlation appearing there.

For two independent XY sectors, let \(F,G\) have supports \(A,A'\), sup-plus-maximum-individual-Lipschitz norms \(L_F,L_G\), and distance at least \(4u\), \(u\geq3\). Let \(A_u\) be the radius-\(u\) neighborhood and let \(\check\partial A_u\) list endpoints of its crossing edges with multiplicity. Then

\begin{equation}
|\operatorname{Cov}_{XY}(F,G)|
\leq C(1+b)^2L_FL_G
\sum_{t=1}^2\left[
\sum_{i\in A_u,j\in A'_u}
+\sum_{i\in A,j\in\check\partial A_u}
+\sum_{i\in A',j\in\check\partial A'_u}
\right](c^t_{ij})^{1/3}.
\label{eq:pre:11-3}
\end{equation}

Here is the localization proof. Conditional on all \(\gamma\)'s, apply \eqref{eq:pre:11-1} to the product sign system for the conditional covariance. Average and use \eqref{eq:pre:11-2}; its cross-support sum is included in the first term of \eqref{eq:pre:11-3}.

For clarity, let \(f(\gamma)=\mathbb E[F\mid\gamma]\) and
\(g(\gamma)=\mathbb E[G\mid\gamma]\), and let \(f_0,g_0\) be the same
conditional expectations after cutting the sign bonds crossing their
respective \(u\)-neighborhoods. We keep the original, full marginal law of
\(\gamma\) throughout this comparison. The cut sign systems factor across
the cuts, so \(f_0\) depends only on \(\gamma|_{A_u}\), and similarly for
\(g_0\). Their suprema remain bounded by those of the original observables.
Differentiation in \(\gamma_i\) gives
\[
 \partial_{\gamma_i}f_0
 =\mathbb E_0[\partial_{\gamma_i}F]
  +\sum_{e\ni i,h}(\partial_{\gamma_i}J_e^h)
                 \Cov_0(F,\tau^h_{e^-}\tau^h_{e^+}).
\]
There are at most eight terms, each coupling derivative is bounded by \(b\),
and the covariance is bounded by \(2\|F\|_\infty\). Thus the individual
Lipschitz constants of \(f_0\) are at most \(C(1+b)L_F\), uniformly in the
neighborhood size. For a merely Lipschitz observable the formula holds
almost everywhere, or follows by smooth approximation with the same
bound. Apply \eqref{eq:pre:11-1} and \eqref{eq:pre:11-2} to \(f_0,g_0\).

Finally control the replacement errors. Interpolate each deleted sign bond. The derivative of a conditional mean is its covariance with the sign-edge product. Apply \eqref{eq:pre:11-1} in the sign variables: it is bounded by
\(CL_F\sum_{i\in A,j\text{ endpoint}}\langle\tau_i\tau_j\rangle\).
Monotonicity bounds every weakened-system correlation by the full conditional sign correlation. Averaging under the original full angle law and using \eqref{eq:pre:11-2} gives the boundary terms of \eqref{eq:pre:11-3}, with a factor \(b\). If \(f,f_0,g,g_0\) are bounded conditional means, the covariance replacement cost is at most \(2\|g\|_\infty\mathbb E|f-f_0|+2\|f_0\|_\infty\mathbb E|g-g_0|\). Thus \(L^1\), not pointwise, control suffices.

\subsubsection{The O(4) decomposition and the second localization}\label{pre:part-11-3}

Write

\[
q_i=(\cos\alpha_i\,z_i,\sin\alpha_i\,w_i),\qquad
\alpha_i\in[0,\pi/2],\quad z_i,w_i\in S^1.
\]

The a priori law is the product of the two uniform circle laws and
\(\sin(2\alpha)\,d\alpha\). Conditional on \(\alpha\), the two circle
sectors are independent XY models, with couplings
\[
 J_e^1=b_e\cos\alpha_x\cos\alpha_y,\qquad
 J_e^2=b_e\sin\alpha_x\sin\alpha_y.
\]
The \(\alpha\) density relative to its product one-site measure is the
product of these two XY partition functions. It is associated by the same
mixed-derivative argument, now using the XY inequalities. The first
derivatives of the cosine-sector couplings are all nonpositive and those
of the sine-sector couplings all nonnegative; their products therefore
have the required sign in each sector. Mixed endpoint derivatives of a
single coupling are nonnegative.

The normalized circle components are bounded measurable one-site functions of \(q\), defined arbitrarily on a Haar-null set, with zero Haar means. The angle \(\alpha-\pi/4\) also has zero Haar mean. Equation \eqref{eq:pre:10-4} therefore gives

\begin{equation}
0\leq\mathbb E_\alpha c^t_{ij}\leq Ce^{-cd(i,j)/s},\qquad
|\operatorname{Cov}(\alpha_i,\alpha_j)|\leq Ce^{-cd(i,j)/s}.
\label{eq:pre:11-4}
\end{equation}

For original local Lipschitz observables \(F,G\) at separation \(d\), use \eqref{eq:pre:11-3} conditionally with \(u=\lfloor d/10\rfloor\). The individual circle Lipschitz norms do not increase, since the map from a circle angle to the corresponding sphere spin has speed at most one. Average and use Jensen,
\(\mathbb E(c^t_{ij})^{1/3}\leq(\mathbb E c^t_{ij})^{1/3}\).
This bounds the averaged conditional covariance by a polynomial times \(e^{-c'd/s}\).

For the remaining covariance of the \(\alpha\)-conditional means, localize again: delete the conditional XY bonds crossing \(A_u\) for the first mean and \(A'_u\) for the second. Keep the original full marginal law of \(\alpha\) fixed. The cut expectations depend only on the corresponding local angle sets. Their individual \(\alpha\)-Lipschitz constants are at most \(C(1+b)L_F\) and \(C(1+b)L_G\): the differentiation formula just displayed now has the cosine of an XY edge in place of a sign-edge product, still bounded by one. Their covariance is bounded by association and \eqref{eq:pre:11-4}.

For their \(L^1\) errors, interpolate a deleted XY bond. Its derivative is a conditional XY covariance of \(F\) with a cosine edge at distance at least \(u-O(1)\). Apply \eqref{eq:pre:11-3} to this pair with radius, for example, \(\lfloor u/10\rfloor\). Every conditional XY two-point function in that inequality is bounded by the full conditional one by coupling monotonicity. Jensen and \eqref{eq:pre:11-4} bound the averaged error. This is the second localization; no estimate for an extended conditional mean is being inferred directly from its original support.

All needed sums have polynomial size: \(|A_u|\leq C|A|(1+d)^2\), the crossing-edge endpoint lists satisfy the same bound, \eqref{eq:pre:11-3} uses pairs of such sets, and the second localization introduces only one additional crossing-edge sum. Small \(d\) uses the trivial bound. For instance, increasing constants gives the convenient common bound

\begin{equation}
\boxed{\quad
|\operatorname{Cov}(F,G)|
\leq C(1+B+|A|+|A'|+d)^{20}L_FL_G e^{-cd/s}.
\quad}
\label{eq:pre:11-5}
\end{equation}

It holds uniformly on the stated unpinned tori and on free subgraphs, for every set of couplings \(0\leq b_e\leq b\), including all intermediate bond weakenings. The support and derivative constants do not grow exponentially with support size.


\subsection{Periodic limits and the partition-function bound}\label{pre:conclusion}

Fix the absolute \(D\) required by the profile parameter dimension, then all other constants, and take \(B=\max(B_*,b)\). The polynomial overhead in \(s=R_*M_*\) is absorbed by \eqref{eq:pre:8-2}, so one may use

\begin{equation}
s(b)\leq\exp[\pi b+C\sqrt{B\log B}].
\label{eq:pre:12-1}
\end{equation}

\subsubsection{Local periodic limits}\label{pre:part-12-1}

For a fixed local Lipschitz observable \(F\), choose a square containing its support at distance \(r\) from its boundary. In any sufficiently large torus, interpolate all bonds crossing that square to zero. The derivative of its expectation is a sum of covariances with those edge-energy observables. Equation \eqref{eq:pre:11-5}, valid uniformly throughout the interpolation, bounds the total change by a polynomial in \(b,r,|\operatorname{supp}F|\) times \(e^{-cr/s(b)}\). After the cut, the law inside the square is the same free-square law for every exterior torus. Letting \(r\to\infty\) proves that every periodic subsequential limit has the same expectation of every local Lipschitz observable. Such functions are dense in the continuous local algebra on the compact spin space. This proves uniqueness of the periodic local limit. It does not claim uniqueness among every possible nonperiodic Gibbs state without another argument.

\subsubsection{Full fixed-lattice spectrum}\label{pre:part-12-2}

In homogeneous even periodic systems, site and link reflection positivity follow from factorization at a site seam and the positive kernel
\(e^{bq\cdot q'}=\sum_{m\geq0}b^m(q\cdot q')^m/m!\) at a link seam. Each
power is positive because it is the inner product of the corresponding
tensor powers of \(q,q'\). These properties first pass to the preceding local limit for continuous
local functions. To extend them to bounded measurable local functions,
approximate each observable in $L^2$ of the marginal on its original
finite support, retaining that support. The finite-spin product is compact
metric, so continuous (and Lipschitz) functions are dense there. Reflection
invariance gives convergence of the reflected copies, and Cauchy--Schwarz
gives convergence of the reflected products. Thus positivity extends
to the bounded measurable local algebra. Define the site-reflection form on bounded local functions in the
nonnegative time half-plane by
\(\langle F,G\rangle=\mathbb E[(\Theta\overline F)G]\), quotient its null
space, and complete. One-step translation \(T\) is symmetric on this dense
space by reflection and translation invariance. For
\(a_n=\|T^nF\|^2\), Schwarz gives \(a_n^2\leq a_{n-1}a_{n+1}\), while
\(a_n\leq\|F\|_\infty^2\). If \(a_1>a_0>0\), log convexity would make
\(a_n\) grow geometrically, a contradiction. If \(a_0=0\), the same
inequality gives \(a_1=0\). Hence \(T\) is well defined on the quotient
and extends to a self-adjoint contraction. Link reflection positivity
gives \(\langle F,TF\rangle\geq0\), so this contraction is nonnegative.

\begin{proof}[Proof of Corollary~\ref{cor:all-temperature}]
For a centered local Lipschitz vector \(F\), its \(n\)-step matrix element is a covariance of two separated local functions. Equation \eqref{eq:pre:11-5} implies
\(\limsup_{n\to\infty}|\langle F,T^nF\rangle|^{1/n}\leq e^{-c/s(b)}\).
Positivity of the spectral measure then places its support in \([0,e^{-c/s(b)}]\). To justify density even when bounded measurable local functions generate
that Hilbert space, fix a finite support. Lipschitz functions are dense in
$L^2$ of its marginal probability measure, since the finite spin product
is a compact metric space. Reflection invariance and Cauchy--Schwarz give
\[
 \|[F]\|_{\mathrm{OS}}^2
 =\mathbb E[(\Theta\overline F)F]\le\mathbb E|F|^2.
\]
Thus $L^2$ approximation implies approximation in the reflection
seminorm; subtracting expectations preserves density in the vacuum
complement. Bounded spectral projections then extend the spectral
inclusion from local Lipschitz vectors to that whole complement.
The periodic-limit uniqueness was proved in the preceding subsection,
and \eqref{eq:pre:12-1} gives the stated quantitative bound.
\end{proof}

This passage explains why neutral-observable mixing is necessary:
it controls the dense local algebra, not only the spin sector.

For large $b$, this preliminary gap is of order at least
$e^{-\pi b-C\sqrt{b\log b}}$. It is insufficient by itself for a positive mass bound in the conventional physical scaling, and no such implication is made here.

\subsubsection{Comparing bond expectations on the two tori}\label{pre:part-12-3}

Let \(e_t(n)=\mathbb E_{t;n,n}(q_0\cdot q_{e_1})\), and fix \(b\) while letting \(0\leq t\leq b\). Use \eqref{eq:pre:11-5} uniformly with the maximum coupling parameter \(b\), not with a varying threshold. For even \(n\geq Cs(b)\), choose the same free square of radius comparable to \(n/4\) around this bond in each of the \(n\)- and \(2n\)-tori. Cut its crossing bonds by interpolation. There are \(O(n)\) such edges; their distance from the fixed bond is \(\geq cn\), their coefficients are at most \(b\), and the edge observables have uniformly bounded sup and individual Lipschitz norms. Therefore

\begin{equation}
|e_t(n)-e_t(2n)|
\leq C(1+B+n)^P e^{-c n/s(b)}
\quad\text{uniformly for }0\leq t\leq b.
\label{eq:pre:12-2}
\end{equation}

After both cuts the central bond has the identical free-square expectation. This comparison does not invoke mixing with pinned boundary fields; every system encountered is an unpinned ferromagnet with some weakened or zero bonds.

Differentiate the actual normalized partition functions. There are \(2n^2\) bonds on the \(n\)-square torus and \(8n^2\) on the \(2n\)-square torus, with equal horizontal and vertical means. Thus

\begin{equation}
\frac{d}{dt}\Delta_t(n)=8n^2[e_t(n)-e_t(2n)],\qquad \Delta_0(n)=0.
\label{eq:pre:12-3}
\end{equation}

Integrating \eqref{eq:pre:12-2} proves the upper bound in \eqref{eq:pre:theorem-delta}, after absorbing \(b\), \(n^2\), and the polynomial into \(C(1+b)^pn^p\). To verify nonnegativity directly, on a row of width \(w\) let
\[
 T_w(u,v)=\exp\left\{\frac b2\sum_i u_i\cdot u_{i+1}
                 +b\sum_i u_i\cdot v_i
                 +\frac b2\sum_i v_i\cdot v_{i+1}\right\}.
\]
The positive crossing kernel, multiplied on both sides by the same
positive function, defines a positive compact operator. Its powers of
order at least two are trace class, and direct integration gives
\(Z(n,w)=\operatorname{Tr}T_w^n\). If its nonnegative eigenvalues are
\(\lambda_j\), then
\(\sum_j\lambda_j^{2n}\leq(\sum_j\lambda_j^n)^2\).
Thus \(p(n,w)=2\log Z(n,w)-\log Z(2n,w)\geq0\). Square-lattice rotation
symmetry now gives the exact identity
\(\Delta(n)=2p(n,n)+p(n,2n)\geq0\).

For large \(b\), the right-hand side of \eqref{eq:pre:12-1} is bounded by \(\Xi(b)\) after increasing its fixed constant. The bounded interval \(0\leq b\leq B_*\) is covered by a further enlargement of that constant. Since \(B=\max(B_*,b)\), its occurrence in the polynomial prefactor can likewise be absorbed into a constant times a power of \(1+b\). This completes the proof of Theorem~\ref{thm:preliminary}.


